% templates/bifurcation_spectrum.tex.mustache
% AUTO-GENERATED Stage 5 bifurcation panel
% Rendered by scripts/build_campaign_report.jl
% DO NOT EDIT MANUALLY

\section{Physical Interpretation of the Continuation Parameter}

The continuation parameter $\gamma$ represents the effective atmosphere--surface coupling strength inferred from the sparse structural operator identified in Stage~4.
Small values of $\gamma$ correspond to weak turbulent transport and strongly decoupled nocturnal conditions, where the surface inversion evolves in relative isolation from the overlying residual layer.
Larger values of $\gamma$ represent enhanced vertical exchange between the surface inversion and the residual layer above, sustaining coherent turbulent structures through the column.

Parameter continuation therefore traces the geometric deformation of the reconstructed attractor manifold as the atmospheric state transitions between coupled and decoupled regimes.
The leading eigenvalue $\alpha(\gamma) = \max\,\Re(\lambda_i)$ serves as a scalar indicator of manifold stability: negative values indicate a contracting, stable attractor; the zero crossing at $\gamma_c$ marks the onset of structural instability; and the imaginary part $\beta(\gamma) = |\mathrm{Im}(\lambda_\mathrm{dom})|$ tracks the intrinsic oscillatory frequency of the emerging mode.

\section{Stage 5 Bifurcation Spectrum}

This panel summarizes the continuation branch spectral abscissa $\alpha(\gamma)$ for campaign \texttt{ GABLS3 }.
The inferred Hopf crossing occurs at $\gamma_c \approx 0.5405$ with local transversality $d\alpha/d\gamma \approx n/a$ and characteristic period $T_H \approx 0.6\,\mathrm{min}$ (35.3 s).

\begin{figure}[htbp]
    \centering
    \usepgfplotslibrary{groupplots}
    \tikzsetnextfilename{bifurcation-panel-gabls3}
    \begin{tikzpicture}
        \begin{groupplot}[
            group style={
                group name=bifurcation_plots,
                group size=1 by 2,
                vertical sep=12pt,
                x descriptions at=edge bottom
            },
            width=0.96\textwidth,
            grid=both,
            grid style={line width=.1pt, draw=gray!15},
            major grid style={line width=.2pt, draw=gray!35},
            xmin=0.402931, xmax=1.000000,
            tick label style={font=\small},
            label style={font=\small},
            axis background/.style={fill=gray!2}
        ]

        \nextgroupplot[
            height=4.5cm,
            ylabel={Growth $\alpha = \max\,\Re(\lambda)$},
            axis lines=middle,
            axis line style={draw=gray!60},
            legend style={at={(0.98,0.05)}, anchor=south east, font=\small},
            xticklabels={}
        ]
            \addplot[
                color=blue!70!black,
                line width=1.3pt,
                mark=none
            ] table[
                col sep=comma,
                x=gamma,
                y=max_real_eig
            ] {../../data/outputs/stage5_bifurcation_branches_gabls3.csv};
            \addlegendentry{$\alpha(\gamma)$}

            \addplot[
                color=red!60,
                dashed,
                line width=1.0pt,
                domain=0.490476:0.590476,
                samples=2
            ] {(x - 0.540476) * (0.0)};
            \addlegendentry{Slope $\left.\frac{d\alpha}{d\gamma}\right|_{\gamma_c}$}

            \addplot[
                color=red,
                mark=*,
                mark size=2.8pt
            ] coordinates {(0.540476, 0.0)};
            \addlegendentry{Hopf crossing}

            \node[anchor=south west, text=red!80!black, font=\small\bfseries] at (axis cs:0.540476, 0.05)
                {$\gamma_c = 0.5405$};

        \nextgroupplot[
            height=3.8cm,
            xlabel={Dissipation parameter ($\gamma$)},
            ylabel={Frequency $\beta\,(\mathrm{rad\,s^{-1}})$},
            ymin=0.0, ymax=0.199360,
            axis lines=left,
            axis line style={draw=gray!60},
            legend style={at={(0.98,0.92)}, anchor=north east, font=\small}
        ]
            \addplot[
                color=teal!80!black,
                line width=1.3pt,
                mark=none
            ] table[
                col sep=comma,
                x=gamma,
                y=max_imag_eig
            ] {../../data/outputs/stage5_bifurcation_branches_gabls3.csv};
            \addlegendentry{$\beta(\gamma)$ frequency branch}

        \end{groupplot}
    \end{tikzpicture}
    \caption{Stage 5 parameter continuation spectrum: the top panel charts the principal manifold stability branch $\alpha(\gamma)$, while the bottom panel shows the aligned oscillation-frequency branch $\beta(\gamma)$ when available.}
    \label{fig:bifurcation_spectrum_gabls3}
\end{figure}


\subsection{Stage~5 Manifold Stability Diagnostics}

\begin{table}[htbp]
\centering
\caption{Stage~5 manifold stability diagnostics for campaign \texttt{ GABLS3 }.
    Values derived from automated parameter continuation;
    \textit{n/a} indicates a diagnostic not resolved in the current sweep.}
\label{tab:stage5_stability_gabls3}
\begin{tabular}{llcc}
\toprule
Diagnostic & Symbol & Value & Units \\
\midrule
Critical coupling parameter & $\gamma_c$ & 0.5405 & -- \\
Transversality slope & $\left.d\alpha/d\gamma\right|_{\gamma_c}$ & n/a & s$^{-1}$ \\
Hopf frequency & $\beta_H$ & 0.17818 & rad\,s$^{-1}$ \\
Oscillation period & $T_H$ & 0.6 & min \\
\bottomrule
\end{tabular}
\end{table}

\subsection{Classical Boundary-Layer Interpretation}

The manifold transition identified at $\gamma_c \approx 0.5405$ coincides with the onset of atmospheric decoupling commonly observed in strongly stable nocturnal boundary layers.
Below the threshold, turbulent transport maintains communication between the surface inversion and the overlying residual flow, and the reconstructed attractor contracts toward a stable equilibrium.
Near $\gamma_c$, the attractor develops a loss of structural stability, indicated by the sign change of the leading eigenvalue, making the boundary layer increasingly susceptible to intermittent turbulence bursts and gravity-wave excitation.
Beyond the threshold, the manifold evolves toward oscillatory dynamics consistent with observed low-level jet pulsations, wave-modulated transport, and the periodic recovery cycles characteristic of very stable conditions.

This geometric framework complements local Richardson-number diagnostics by providing a globally consistent indicator of the boundary-layer stability state that remains valid through and beyond the classical collapse threshold of $\text{Ri}_g = 0.25$.


\noindent\textit{Cross-campaign Stage 5 comparison table omitted because one or more required summary manifests were not available at render time.}
